\chapter{Urinary Retention}

\section*{Definition}
Inability to partially or completely empty the bladder, classified as acute (sudden and painful) or chronic (gradual onset with less discomfort but significant residual urine volume).

\section*{Pathophysiology}
Acute retention typically results from mechanical obstruction (BPH, urethral stricture, stone, clot) or neurological causes (spinal cord injury, post-operative). Chronic retention involves detrusor underactivity or bladder outlet obstruction with compensatory detrusor hypertrophy, leading to high post-void residuals, overflow incontinence, and upper tract damage risk.

\section*{Etiology}
\begin{itemize}
  \item Benign prostatic hyperplasia (BPH) — most common cause in men
  \item Urethral stricture
  \item Medication side effects (anticholinergics, antihistamines, decongestants)
  \item Post-operative (anesthesia, opioids, pelvic surgery)
  \item Neurological conditions (spinal cord injury, multiple sclerosis, diabetic neuropathy)
  \item Severe constipation (fecal impaction)
  \item Bladder stones or tumors
  \item Pelvic organ prolapse (cystocele) in women
  \item Prostate cancer or bladder cancer
  \item Phimosis or paraphimosis
\end{itemize}

\section*{Indications for Evaluation}
Seek care if:
\begin{itemize}
  \item Severe or worsening symptoms
  \item Complete inability to urinate (emergency requiring catheterization), painful distension, or acute retention with fever
  \item Accompanied by other concerning signs
\end{itemize}

\section*{Related Symptoms}
\begin{itemize}
  \item Urinary frequency
  \item Hesitancy
  \item Weak urinary stream
  \item Overflow incontinence
  \item Abdominal distension
\end{itemize}
\textbf{Note:} Always consider serious underlying conditions when evaluating persistent urinary retention.

\section*{Self-Care / Home Management}
\begin{itemize}
  \item Rest and avoid activities that may aggravate the condition.
  \item Maintain adequate hydration and a balanced diet.
  \item Over-the-counter remedies may provide symptomatic relief (consult a pharmacist or doctor).
  \item Monitor symptoms and seek professional advice if they persist or worsen.
  \item Avoid self-medicating with prescription medications without medical guidance.
\end{itemize}

\section*{Diagnostic Approach}
\begin{itemize}
  \item A thorough history and physical examination are the first steps in evaluation.
  \item Laboratory tests (blood work, urinalysis) may be ordered based on clinical suspicion.
  \item Imaging studies (X-ray, ultrasound, CT, MRI) may be indicated when structural pathology is suspected.
  \item Specialist referral is arranged when the diagnosis remains uncertain or requires specific expertise.
\end{itemize}

\section*{Treatment / Management Overview}
\begin{itemize}
  \item Treatment targets the underlying cause identified through diagnostic evaluation.
  \item Supportive care and symptom management are provided while awaiting definitive therapy.
  \item Medications, lifestyle modifications, and physical therapies may be prescribed as appropriate.
  \item Follow-up ensures treatment effectiveness and allows timely adjustments to the care plan.
\end{itemize}

\clearpage